\documentclass[aps,prd,preprint,nofootinbib,superscriptaddress]{revtex4-2}

\usepackage[T1]{fontenc}
\usepackage[utf8]{inputenc}
\usepackage{amsmath,amssymb,mathtools}
\usepackage{bm}
\usepackage{booktabs}
\usepackage{graphicx}
\usepackage{xcolor}
\usepackage[colorlinks=true,linkcolor=blue,citecolor=blue,urlcolor=blue]{hyperref}

\newcommand{\ii}{\mathrm{i}}
\newcommand{\Ds}{D_s}
\newcommand{\Dst}{D_s^*}
\newcommand{\Dsone}{D_{s1}}
\newcommand{\Bs}{B_s}
\newcommand{\Bst}{B_s^*}
\newcommand{\Bsone}{B_{s1}}
\newcommand{\GeV}{\mathrm{GeV}}
\newcommand{\MeV}{\mathrm{MeV}}
\newcommand{\keV}{\mathrm{keV}}
\newcommand{\qqbar}{\langle\bar q q\rangle}
\newcommand{\ssbar}{\langle\bar s s\rangle}

\begin{document}



\title{Radiative transitions of strange heavy axial mesons in light-cone QCD sum rules}
%\title{Spin mixing and radiative transitions of strange heavy axial mesons in light-cone QCD sum rules}

\author{S. Bilmis}
\affiliation{Department of Physics, Middle East Technical University, Ankara, Turkey}
\author{T. M. Aliev}
\affiliation{Department of Physics, Middle East Technical University, Ankara, Turkey}
\author{M. Savci}
\affiliation{Department of Physics, Middle East Technical University, Ankara, Turkey}

\date{\today}

\begin{abstract}
We calculate the radiative widths of
$\Dsone(2460)\to\Ds\gamma$, $\Dsone(2536)\to\Ds\gamma$,
$\Bsone(5750)\to\Bs\gamma$, and $\Bsone(5830)\to\Bs\gamma$
in light-cone QCD sum rules.  The two basis-current sum rules are combined
into the physical axial-vector amplitudes using the same mixing convention
as in the hadronic residues.  With the lattice-normalized leading-twist photon
input, we obtain
\[
\Gamma[\Dsone(2460)\to\Ds\gamma]
   =32.7^{+6.7}_{-6.2}~\keV,
\qquad
\Gamma[\Dsone(2536)\to\Ds\gamma]
   =0.425^{+1.352}_{-0.390}~\keV .
\]
The smallness of the $\Dsone(2536)\to\Ds\gamma$ width is traced to destructive
interference between the two physical-current components.  In the
bottom-strange sector we find
\[
\Gamma[\Bsone(5750)\to\Bs\gamma]
   =19.3^{+7.9}_{-5.2}~\keV,
\qquad
\Gamma[\Bsone(5830)\to\Bs\gamma]
   =0.504^{+1.173}_{-0.447}~\keV ,
\]
where the lower state is treated as a theoretical diagnostic with
$M_{\Bsone(5750)}\simeq5.75~\GeV$.  The results support a strong hierarchy
between the radiative widths of the two axial states and provide predictions
for future measurements of bottom-strange radiative transitions.
\end{abstract}

\maketitle



\section{Introduction}
\label{sec:intro}

Radiative transitions of heavy-light mesons provide a sensitive probe of their
internal structure.  In positive-parity strange heavy mesons, such transitions
are especially useful because they depend on the spin composition of the
physical axial-vector states and on the coupling of the photon to both the
heavy and strange quark degrees of freedom.  They therefore complement mass
spectroscopy and strong-decay measurements in testing the assignment of the
observed states and the role of spin mixing.

The charm-strange axial-vector mesons $\Dsone(2460)$ and $\Dsone(2536)$ are
particularly important examples.  The $\Dsone(2460)$ is a very narrow state
below the $D^{\ast}K$ threshold and its radiative decay to $\Ds\gamma$ has
been observed in $B$ decays and in continuum production
\cite{BaBar:2004zjt,BaBar:2006eep}.  The $\Dsone(2536)$, on the other hand,
is well established in hadronic channels such as $D^{\ast}K$ and
$D\pi K$ \cite{Belle:2007hue,BaBar:2011zpy}, but its radiative decay to
$\Ds\gamma$ has not been observed with comparable significance.  The measured
width of the $\Dsone(2536)$ is below the MeV scale, which makes a suppressed
radiative branching fraction experimentally difficult to isolate
\cite{BaBar:2011zpy}.  In the bottom-strange sector the
$\Bsone(5830)$ state has been observed in $B^{\ast}K$ spectroscopy
\cite{LHCb:2012kxf,CMS:2018qdl}, whereas no direct information is available
on its radiative transition to $\Bs\gamma$.

Theoretical studies of these radiative modes have been carried out using
constituent quark models, relativistic quark models, effective approaches,
and QCD sum rules
\cite{Goity:2000dk,Godfrey:2005ww,Colangelo:2005hv,Radford:2009bs,
Godfrey:2016nwn,Lu:2016bbk,Green:2016occ,Chen:2020ejk,Li:2021qgz,
Pullin:2021ebn,Bondar:2023mua,Bondar:2025smw}.  The light-cone QCD sum-rule
calculation of Ref.~\cite{Colangelo:2005hv} showed that
$\Dsone(2460)\to\Ds\gamma$ is expected to be one of the leading radiative
channels of the $\Dsone(2460)$.  More recently, Bondar and Milstein argued
that the smallness of $\Dsone(2536)\to\Ds\gamma$ follows from a cancellation
between the two spin components of the physical axial-vector state
\cite{Bondar:2023mua,Bondar:2025smw}.  Their analysis also considered the
corresponding bottom-strange transitions, including both the established
$\Bsone(5830)$ state and a lower \(\Bsone(5750)\) configuration.
The spread among existing predictions is particularly large for the channels
where the two axial-current components interfere destructively.

In this work we calculate the radiative transitions
$\Dsone(2460)\to\Ds\gamma$, $\Dsone(2536)\to\Ds\gamma$,
$\Bsone(5750)\to\Bs\gamma$, and $\Bsone(5830)\to\Bs\gamma$
in light-cone QCD sum rules.  We keep the strange-quark-mass terms and use
photon distribution amplitudes, including the recent lattice-normalized input
for the leading-twist photon distribution amplitude.  The two basis-current
sum rules are combined into the physical axial-vector amplitudes, which allows
us to test whether the cancellation mechanism proposed for
$\Dsone(2536)\to\Ds\gamma$ is reproduced in the sum-rule framework.  For the
bottom-strange sector, the lower \(\Bsone(5750)\) state is treated as a
theoretical diagnostic, while the $\Bsone(5830)$ result is a prediction for
the experimentally established state.

The paper is organized as follows. In Sec.~\ref{sec:framework} we define the
interpolating currents, the physical-current mixing, the radiative form factor,
and the light-cone sum rules. Section~\ref{sec:numerics} contains the numerical
analysis, comparison with previous studies, and discussion of the experimental
status. We summarize our results in Sec.~\ref{sec:concl}.

\section{Theoretical framework}
\label{sec:framework}

We consider the radiative transition of an axial-vector strange heavy meson
with momentum \(p^{\prime}=p+q\) into a pseudoscalar strange heavy meson with momentum
\(p\) and a real photon with momentum \(q\).  The heavy quark is denoted by
\(Q=c,b\), while the light quark is \(q=s\) in the applications considered
below.  The pseudoscalar interpolating current is chosen as
\begin{equation}
  J_5=\bar q\,\ii\gamma_5 Q .
  \label{eq:j5}
\end{equation}
For the axial-vector channel we use two independent basis currents,
\begin{align}
  J^A_\mu &= \bar q\gamma_\mu\gamma_5 Q, \\
  J^B_\mu &= \frac{\ii}{m_Q+m_q}\,
  \bar q\sigma_{\mu\alpha}p^{\prime\alpha}\gamma_5 Q ,
  \label{eq:basis-currents}
\end{align}
where \(p^{\prime\alpha}\) is the momentum of the initial axial-vector meson.  The
physical-current combinations are written as
\begin{align}
  J_{1\mu} &= \sin\theta\,J^A_\mu+\cos\theta\,J^B_\mu,\\
  J_{2\mu} &= \cos\theta\,J^A_\mu-\sin\theta\,J^B_\mu .
  \label{eq:physical-currents}
\end{align}
With this convention, \(J_{1\mu}\) is assigned to the lower axial-vector
combination, whereas \(J_{2\mu}\) is assigned to the higher one.  Thus, in the
charm-strange sector, \(J_1\) is associated with \(\Dsone(2460)\) and \(J_2\)
with \(\Dsone(2536)\).  In the bottom-strange sector, the same convention
associates \(J_1\) with the lower diagnostic configuration \(\Bsone(5750)\)
and \(J_2\) with the established \(\Bsone(5830)\).  We use
\(\theta=35.3^\circ\) as the central value and vary it in the interval
\(25^\circ\leq\theta\leq45^\circ\) in the numerical analysis.

The decay constants are defined by
\begin{align}
  \langle 0|\bar q\,\ii\gamma_5 Q|P(p)\rangle
  &=
  \frac{m_P^2 f_P}{m_Q+m_q},
  \label{eq:fp-def}\\
  \langle0|J_{i\mu}|X_i(p^{\prime},\eta)\rangle
  &=f_i m_i\eta_\mu ,
  \qquad i=1,2 ,
  \label{eq:fi-def}
\end{align}
where \(X_i\) denotes the axial-vector state coupled to \(J_{i\mu}\), with
mass \(m_i\), polarization vector \(\eta_\mu\), and decay constant \(f_i\).

The radiative matrix element is parametrized by a single gauge-invariant
electric-dipole form factor,
\begin{equation}
  \langle P(p)\gamma(q,\varepsilon)|X_i(p^{\prime},\eta)\rangle
  =
  e g_i
  \left[
  (\eta\cdot\varepsilon^*)(p\cdot q)
  -(\eta\cdot q)(p\cdot\varepsilon^*)
  \right].
  \label{eq:e1}
\end{equation}
The corresponding partial width is
\begin{equation}
  \Gamma(X_i\to P\gamma)
  =
  \frac{\alpha_{\rm em}}{3}\,g_i^2
  \left(
  \frac{m_i^2-m_P^2}{2m_i}
  \right)^3 .
  \label{eq:width}
\end{equation}

Following the three-point notation of Rohrwild~\cite{Rohrwild:2007yt}, the
light-cone sum rule starts from
\begin{equation}
  \Pi_{i\mu\nu}(p,q)=\ii^2\int d^4x\,d^4y\,
  e^{\ii p\cdot x+\ii q\cdot y}
  \langle0|T\{J_5(x)j_\nu^{\rm em}(y)J_{i\mu}^{\dagger}(0)\}|0\rangle,
  \qquad i=1,2,
  \label{eq:corr-physical}
\end{equation}
where
\begin{equation}
  j_\nu^{\rm em}=e_q\bar q\gamma_\nu q+e_Q\bar Q\gamma_\nu Q
\end{equation}
and the charges are in units of the positive elementary charge.  Equivalently,
introduce the plane-wave background field
\begin{equation}
  F_{\alpha\beta}(z)=\ii
  (q_\alpha\varepsilon_\beta^*-q_\beta\varepsilon_\alpha^*)
  e^{\ii q\cdot z}.
\end{equation}
The term linear in $F_{\alpha\beta}$ satisfies
\begin{equation}
  \varepsilon^{*\nu}\Pi_{i\mu\nu}(p,q)
  =\ii\int d^4x\,e^{\ii p\cdot x}
  \langle0|T\{J_5(x)J_{i\mu}^{\dagger}(0)\}|0\rangle_F.
  \label{eq:corr-background}
\end{equation}
The commonly used vacuum-to-photon correlator is precisely this linear
background-field component, rather than an independent starting convention.
Since \(J_{1\mu}\) and \(J_{2\mu}\) are linear combinations of \(J^A_\mu\) and
\(J^B_\mu\), we first calculate the two basis correlators and only afterwards
form the physical combinations.  We use the coefficient of the gauge-invariant
structure
\begin{equation}
  (p\cdot q)\varepsilon_\mu^*
  -(p\cdot\varepsilon^*)q_\mu .
  \label{eq:projector}
\end{equation}
For the basis currents we write
\begin{equation}
  \Pi^K_\mu(p,q)=
  \Pi_K(p^2,p^{\prime2})
  \left[
  (p\cdot q)\varepsilon_\mu^*
  -(p\cdot\varepsilon^*)q_\mu
  \right]
  +\cdots ,
  \qquad K=A,B ,
  \label{eq:basis-decomp}
\end{equation}
where the remaining Lorentz structures are not used in the sum rule.

On the hadronic side, isolating the ground-state double pole gives
\begin{equation}
  \Pi^{\rm had}_{i}(p^2,p^{\prime2})
  =
  \frac{m_P^2 f_P}{m_Q+m_q}\,
  \frac{m_i f_i g_i}
  {(m_i^2-p^{\prime2})(m_P^2-p^2)}
  +\cdots ,
  \qquad i=1,2 ,
  \label{eq:hadronic-invariant}
\end{equation}
where the ellipsis denotes higher-state and continuum contributions.

The QCD side is obtained from the light-cone expansion in the external
electromagnetic field.  For the heavy-quark line we use
\begin{align}
  S_Q(k) &=
  \frac{\not k+m_Q}{m_Q^2-k^2}
  -e_Q\int_0^1du
  \left[
  \frac{\not k+m_Q}{2(m_Q^2-k^2)^2}
  F^{\mu\nu}(ux)\sigma_{\mu\nu}
  +\frac{u x_\mu F^{\mu\nu}(ux)\gamma_\nu}{m_Q^2-k^2}
  \right]
  \notag\\
  &\quad
  -g_s\int_0^1du
  \left[
  \frac{\not k+m_Q}{2(m_Q^2-k^2)^2}
  G^{\mu\nu}(ux)\sigma_{\mu\nu}
  +\frac{u x_\mu G^{\mu\nu}(ux)\gamma_\nu}{m_Q^2-k^2}
  \right]
  +\cdots .
  \label{eq:heavy-prop}
\end{align}
The first term gives the free heavy-quark propagator.  The terms proportional
to \(e_Q\) describe perturbative photon emission from the heavy-quark line,
whereas the terms proportional to \(g_s\) generate the quark-gluon photon
matrix elements.  For the strange line we use Rohrwild's coordinate-space
notation, but retain the mass terms omitted in the massless expression of
Ref.~\cite{Rohrwild:2007yt}.  With $\bar u=1-u$,
\begin{align}
  \widehat{s^a(x)\bar s^b(0)}={}&
  \delta^{ab}\Bigg[
    \frac{\ii\not x}{2\pi^2x^4}-\frac{m_s}{4\pi^2x^2}
  \Bigg]
  \notag\\
  &-\frac{\ii g_s}{16\pi^2x^2}\int_0^1du\,
  \{\bar u\not x\sigma_{\alpha\beta}
  +u\sigma_{\alpha\beta}\not x\}
  G_A^{\alpha\beta}(ux)t^A_{ab}
  \notag\\
  &-\frac{\ii e_s\delta^{ab}}{16\pi^2x^2}\int_0^1du\,
  \{\bar u\not x\sigma_{\alpha\beta}
  +u\sigma_{\alpha\beta}\not x\}F^{\alpha\beta}(ux)
  \notag\\
  &-\frac{m_s}{32\pi^2}\int_0^1du\,
  \left[g_sG_A^{\alpha\beta}(ux)t^A_{ab}
  +e_s\delta^{ab}F^{\alpha\beta}(ux)\right]
  \sigma_{\alpha\beta}L_x+\cdots,
  \label{eq:light-prop-rohrwild}
\end{align}
where
\begin{equation}
  L_x=\ln\left(\frac{-x^2\mu^2}{4}\right)+2\gamma_E.
\end{equation}
No local vacuum-condensate term is displayed in
Eq.~\eqref{eq:light-prop-rohrwild}.  Setting $m_s$ and the last line to zero
recovers Rohrwild's displayed light propagator.  The nonlocal
vacuum-to-photon quark and quark-gluon matrix elements define the two- and
three-particle photon DAs collected in
Appendix~\ref{app:sumrule-expressions}.

We emphasize an implementation point relevant to the numerical tables: the
existing axial-current code still includes the local heavy-line photon
contribution named \texttt{heavy\_local}, proportional to
$e_Q e^{-m_Q^2/M^2}\langle\bar s s\rangle$.  It is not part of the displayed
propagator above and has not been silently removed from the numerical results.
The mixed and vacuum-gluon condensate contributions were not included in the
radiative three-point calculation.

After the double Borel transformation in \(p^{\prime2}\) and \(p^2\) and continuum
subtraction, the basis-current invariant amplitudes are written as
\begin{equation}
  \widehat\Pi_K(M_1^2,M_2^2,s_0)
  =
  \widehat\Pi_K^{\rm pert}
  +\delta_{KA}\widehat\Pi_A^{\rm local}
  +\widehat\Pi_K^{(2p)}
  +\widehat\Pi_K^{(3p)},
  \qquad K=A,B .
  \label{eq:pi-decomp}
\end{equation}
Here the terms denote perturbative photon emission, the separately implemented
local light-quark-condensate contribution in the present axial calculation,
two-particle photon-DA contributions, and three-particle photon-DA
contributions, respectively.  No local condensate term is thereby being
inserted into either propagator displayed above.  At the symmetric Borel point
used in the numerical analysis,
\begin{equation}
  M_1^2=M_2^2=2M^2,
  \qquad
  u_0=\frac{M_1^2}{M_1^2+M_2^2}=\frac12 ,
  \label{eq:borel-defs}
\end{equation}
the perturbative parts have the generic form
\begin{equation}
  \widehat\Pi_K^{\rm pert}
  =
  \int_{(m_Q+m_s)^2}^{s_0}ds\,e^{-s/M^2}\rho_K(s),
  \qquad K=A,B .
  \label{eq:pert-part}
\end{equation}
The photon-DA conventions and auxiliary combinations entering
\(\widehat\Pi_K^{(2p)}\) and \(\widehat\Pi_K^{(3p)}\) are given in
Appendix~\ref{app:sumrule-expressions}.

The physical invariant amplitudes are obtained by applying the same rotation
as in Eq.~\eqref{eq:physical-currents},
\begin{align}
  \widehat\Pi_1
  &=
  \sin\theta\,\widehat\Pi_A
  +\cos\theta\,\widehat\Pi_B ,
  \label{eq:pi1-rotation}\\
  \widehat\Pi_2
  &=
  \cos\theta\,\widehat\Pi_A
  -\sin\theta\,\widehat\Pi_B .
  \label{eq:pi2-rotation}
\end{align}
Matching the hadronic and QCD representations gives
\begin{equation}
  g_i =
  \frac{m_Q+m_q}{m_P^2 f_P\,m_i f_i}
  \exp\left(
  \frac{m_i^2}{M_1^2}
  +\frac{m_P^2}{M_2^2}
  \right)
  \widehat\Pi_i(M_1^2,M_2^2,s_0),
  \qquad i=1,2 .
  \label{eq:general-sum-rule}
\end{equation}
Equivalently,
\begin{align}
  g_1 &=
  \frac{m_Q+m_q}{m_P^2 f_P\,m_1 f_1}
  \exp\left(
  \frac{m_1^2}{M_1^2}
  +\frac{m_P^2}{M_2^2}
  \right)
  \left[
  \sin\theta\,\widehat\Pi_A
  +\cos\theta\,\widehat\Pi_B
  \right],
  \label{eq:g1-sum-rule}\\
  g_2 &=
  \frac{m_Q+m_q}{m_P^2 f_P\,m_2 f_2}
  \exp\left(
  \frac{m_2^2}{M_1^2}
  +\frac{m_P^2}{M_2^2}
  \right)
  \left[
  \cos\theta\,\widehat\Pi_A
  -\sin\theta\,\widehat\Pi_B
  \right].
  \label{eq:g2-sum-rule}
\end{align}
The first expression is used for
\(\Dsone(2460)\to\Ds\gamma\) and
\(\Bsone(5750)\to\Bs\gamma\), while the second is used for
\(\Dsone(2536)\to\Ds\gamma\) and
\(\Bsone(5830)\to\Bs\gamma\).

\subsection{Physical-current two-point functions and decay constants}

For clarity, the physical residues entering Eqs.~\eqref{eq:g1-sum-rule} and
\eqref{eq:g2-sum-rule} are associated with the transverse two-point matrix
\begin{equation}
  \Pi_{\mu\nu}^{XY}(p^{\prime})=
  \ii\int d^4x\,e^{\ii p^{\prime}\cdot x}
  \langle0|T\{J_\mu^X(x)J_\nu^{Y\dagger}(0)\}|0\rangle,
  \qquad X,Y=A,B.
  \label{eq:twopoint-basis}
\end{equation}
Suppressing the common transverse tensor, the same state rotation gives
\begin{align}
  \Pi_{11}&=\sin^2\theta\,\Pi_{AA}+\cos^2\theta\,\Pi_{BB}
  +2\sin\theta\cos\theta\,\Pi_{AB},
  \label{eq:pi11}\\
  \Pi_{22}&=\cos^2\theta\,\Pi_{AA}+\sin^2\theta\,\Pi_{BB}
  -2\sin\theta\cos\theta\,\Pi_{AB},
  \label{eq:pi22}\\
  \Pi_{12}&=\sin\theta\cos\theta(\Pi_{AA}-\Pi_{BB})
  +(\cos^2\theta-\sin^2\theta)\Pi_{AB}.
  \label{eq:pi12}
\end{align}
If
\begin{equation}
  \widehat\Pi_{XY}=\int_{s_{\rm th}}^{s_0}ds\,
  e^{-s/{\cal M}^2}\rho_{XY}^{\rm pert}(s)
  +\sum_d\widehat\Pi_{XY}^{(d)},
  \label{eq:twopoint-ope-blocks}
\end{equation}
the explicit physical combination can be kept compact as
\begin{align}
  \widehat\Pi_{11}={}&\int_{s_{\rm th}}^{s_0}ds\,
  e^{-s/{\cal M}^2}
  \left[\sin^2\theta\rho_{AA}^{\rm pert}
  +\cos^2\theta\rho_{BB}^{\rm pert}
  +2\sin\theta\cos\theta\rho_{AB}^{\rm pert}\right]
  \notag\\
  &+\sum_d\left[\sin^2\theta\widehat\Pi_{AA}^{(d)}
  +\cos^2\theta\widehat\Pi_{BB}^{(d)}
  +2\sin\theta\cos\theta\widehat\Pi_{AB}^{(d)}\right],
  \label{eq:pi11-ope-explicit}
\end{align}
with $\widehat\Pi_{22}$ obtained by the replacements shown explicitly in
Eq.~\eqref{eq:pi22}.  The corresponding direct sum rules would be
\begin{align}
  f_1^2&=\frac{e^{m_1^2/{\cal M}^2}}{m_1^2}\widehat\Pi_{11},&
  f_2^2&=\frac{e^{m_2^2/{\cal M}^2}}{m_2^2}\widehat\Pi_{22}.
  \label{eq:f12-direct}
\end{align}
In the present calculation the $AA$ two-point OPE is available, but the full
$BB$ and $AB$ OPEs have not been independently completed.  Therefore the
numerical $f_1,f_2$ values quoted below are overlap-model normalizations, not
independent predictions of Eqs.~\eqref{eq:twopoint-ope-blocks} and
\eqref{eq:f12-direct}.

\section{Numerical analysis}
\label{sec:numerics}

The external inputs and benchmark decay constants used in the numerical
analysis are collected in Table~\ref{tab:inputs}.  The quark masses and hadron
masses are taken from the Particle Data Group~\cite{ParticleDataGroup:2024cfk}.
The pseudoscalar decay constants are taken from the lattice averages
summarized by FLAG~\cite{FlavourLatticeAveragingGroupFLAG:2024oxs}.  For the
photon distribution amplitudes we use the BBK/Rohrwild convention
\cite{Ball:2002ps,Rohrwild:2007yt}.  We compare two normalizations of the
leading-twist photon matrix element: the conventional susceptibility
normalization, proportional to $\chi_s\langle\bar s s\rangle$, and the recent
lattice normalization $f_{\gamma,s}^{\perp}$ of the leading-twist photon
DA~\cite{Bacchio:2024pij}.  The latter is used for our preferred numerical
predictions, while the susceptibility-normalized results are retained to
facilitate comparison with earlier photon-LCSR analyses.

\begin{table}[h]
\caption{External inputs and benchmark decay constants used in the numerical
analysis.  Photon-DA normalizations are quoted at $\mu=1~\GeV$.}
\label{tab:inputs}
\begin{ruledtabular}
\begin{tabular}{ll}
Parameter & Value \\
\hline
$m_c,\ m_b,\ m_s$
& $1.27(2),\ 4.18(3),\ 0.093(11)~\GeV$
  \cite{ParticleDataGroup:2024cfk} \\
$M_{\Dsone(2460)},\ M_{\Dsone(2536)},\ m_{\Ds}$
& $2.4595,\ 2.53511,\ 1.96835~\GeV$
  \cite{ParticleDataGroup:2024cfk} \\
$M_{\Bsone(5830)},\ m_{\Bs}$
& $5.82870,\ 5.36692~\GeV$
  \cite{ParticleDataGroup:2024cfk,CMS:2018qdl} \\
$M_{\Bsone(5750)}$
& $5.750(26)~\GeV$ \cite{Lang:2015hza} \\
$f_{\Ds},\ f_{\Bs}$
& $0.2499,\ 0.2303~\GeV$
  \cite{FlavourLatticeAveragingGroupFLAG:2024oxs} \\
$f_A(\Dsone)$
& $0.245(17)~\GeV$ \cite{Wang:2015mxa} \\
$f_A^{(B_s)},\ f_T^{(B_s)}$
& $0.305(27),\ 0.285(27)~\GeV$
  \cite{Pullin:2021ebn} \\
$\chi_s(1~\GeV)$
& $+3.15(30)~\GeV^{-2}$ \cite{Rohrwild:2007yt} \\
$f_{\gamma,s}^{\perp}(1~\GeV)$
& $-55.1(1.9)~\MeV$ \cite{Bacchio:2024pij} \\
\end{tabular}
\end{ruledtabular}
\end{table}

The physical-current residues are not independent external inputs.  The
reduced overlap model used in the numerical scan is
\begin{align}
 f_1^2&=\sin^2\theta f_A^2+\cos^2\theta f_B^2
 +2\sin\theta\cos\theta\rho_{AB}f_Af_B,\notag\\
 f_2^2&=\cos^2\theta f_A^2+\sin^2\theta f_B^2
 -2\sin\theta\cos\theta\rho_{AB}f_Af_B.
 \label{eq:residue-overlap-model}
\end{align}
For the charm-strange sector we use $f_1=0.345(17)~\GeV$ as the working
normalization anchor, determine $\rho_{AB}$ from the first equation, and then
obtain $f_2$.  At the central angle the saved scans give
$f_1=0.344[0.329,0.363]~\GeV$ and
$f_2=0.380[0.339,0.423]~\GeV$ for the legacy photon input, and
$f_1=0.345[0.330,0.364]~\GeV$ and
$f_2=0.379[0.338,0.423]~\GeV$ for the lattice-normalized input.  For the
bottom-strange sector, the basis-current constants of
Ref.~\cite{Pullin:2021ebn} fix the basis normalization and $\rho_{AB}$ is
chosen by the closure condition $\Pi_{12}=0$.  For the $\Bsone(5830)$ scan
this gives $f_2=0.218[0.138,0.292]~\GeV$ (legacy) and
$0.203[0.121,0.291]~\GeV$ (lattice).  These are the numerical values actually
found, with the model qualification stated below Eq.~\eqref{eq:f12-direct}.

The Borel windows and continuum thresholds are summarized in
Table~\ref{tab:windows}.  The windows are chosen by the standard sum-rule
criteria: the prediction should be stable against variations of the Borel
parameter, the continuum contribution should not dominate the dispersion
integral, and the higher-twist photon-DA terms should remain under control.
Since the initial and final mesons in each channel are close in mass, we use
the symmetric Borel choice
\begin{equation}
  M_1^2=M_2^2=2M^2,
  \qquad
  u_0=\frac12 ,
  \label{eq:numerical-borel-choice}
\end{equation}
as in Sec.~\ref{sec:framework}.

\begin{table}[t]
\caption{Auxiliary parameters used in the numerical analysis.}
\label{tab:windows}
\begin{ruledtabular}
\begin{tabular}{lll}
Channel & $M^2$ window & $s_0$ window \\
\hline
$\Dsone(2460),\Dsone(2536)\to\Ds\gamma$
& $3.0$--$4.5~\GeV^2$
& $7.5$--$8.5~\GeV^2$ \\
$\Bsone(5750)\to\Bs\gamma$
& $10.0$--$14.0~\GeV^2$
& $36.5$--$40.5~\GeV^2$ \\
$\Bsone(5830)\to\Bs\gamma$
& $10.0$--$14.0~\GeV^2$
& $38.0$--$41.0~\GeV^2$ \\
\end{tabular}
\end{ruledtabular}
\end{table}

The central mixing angle is taken as $\theta=35.3^\circ$, corresponding to the
heavy-quark-limit value for the mixing of the two axial-vector configurations.
To estimate the sensitivity to deviations from this limit and to
phenomenological convention choices, we also scan the interval
$25^\circ\leq\theta\leq45^\circ$.  This scan is treated as a systematic
variation rather than as a statistical error.  The angle dependence is
particularly important for the channels in which the two basis-current
contributions interfere destructively.

Representative stability plots are shown in Fig.~\ref{fig:stability}.  In the
chosen working windows the form factors show acceptable stability under
variations of $M^2$ and $s_0$.  The residual dependence on these auxiliary
parameters is included in the quoted uncertainties.

\begin{figure}[t]
\centering
%\includegraphics[width=0.47\textwidth]{figures/stability_charm.pdf}
%\includegraphics[width=0.47\textwidth]{figures/stability_bottom.pdf}
\caption{Representative Borel and continuum-threshold stability checks for
the charm-strange and bottom-strange channels.}
\label{fig:stability}
\end{figure}

The uncertainties are obtained by varying the inputs listed in
Tables~\ref{tab:inputs} and~\ref{tab:windows}.  Dimensionful quantities with
quoted errors are sampled within their stated uncertainties, while the Borel
parameter and continuum threshold are varied over the working intervals in
Table~\ref{tab:windows}.  Where indicated, the mixing angle is also varied in
the interval $25^\circ\leq\theta\leq45^\circ$.  For each sample, the physical
form factor is computed from Eq.~\eqref{eq:general-sum-rule}, and the
corresponding width is obtained from Eq.~\eqref{eq:width}.  We quote medians
and $16$--$84$ percentile intervals.  This convention naturally gives
asymmetric uncertainties and is more appropriate than a Gaussian fit for the
cancellation-sensitive channels, especially
$\Dsone(2536)\to\Ds\gamma$ and $\Bsone(5830)\to\Bs\gamma$.

The final form factors and widths at the central mixing angle are listed in
Table~\ref{tab:results}.  The $f_{\gamma,s}^{\perp}$ rows are our preferred
predictions.  The $\chi_s\langle\bar s s\rangle$ rows are retained to show the
size of the photon-normalization dependence and to allow comparison with
earlier photon-LCSR calculations.

\begin{widetext}
\begin{table}[t]
\caption{Form factors and radiative widths obtained at the central mixing
angle $\theta=35.3^\circ$.  Here $f_X$ denotes the physical-current residue of
the initial axial-vector state entering Eq.~\eqref{eq:general-sum-rule}.  The
quoted uncertainties correspond to the $16$--$84$ percentile intervals.}
\label{tab:results}
\scriptsize
\begin{ruledtabular}
\begin{tabular}{llccc}
Channel & Photon input & $f_X~(\GeV)$
& $g_X~(\GeV^{-1})$ & $\Gamma~(\keV)$\\
\hline
$\Dsone(2460)\to\Ds\gamma$
& $\chi_s\langle\bar s s\rangle$
& $0.347^{+0.017}_{-0.018}$
& $-0.555^{+0.077}_{-0.075}$
& $64.8^{+18.4}_{-16.8}$\\

$\Dsone(2536)\to\Ds\gamma$
& $\chi_s\langle\bar s s\rangle$
& $0.382^{+0.043}_{-0.043}$
& $+0.0031^{+0.0558}_{-0.0667}$
& $0.640^{+1.13}_{-0.589}$\\

$\Dsone(2460)\to\Ds\gamma$
& $f_{\gamma,s}^{\perp}$
& $0.345^{+0.016}_{-0.016}$
& $-0.394^{+0.039}_{-0.039}$
& $32.7^{+6.7}_{-6.2}$\\

$\Dsone(2536)\to\Ds\gamma$
& $f_{\gamma,s}^{\perp}$
& $0.380^{+0.041}_{-0.043}$
& $+0.0329^{+0.0427}_{-0.0432}$
& $0.425^{+1.352}_{-0.390}$\\

$\Bsone(5750)\to\Bs\gamma$
& $\chi_s\langle\bar s s\rangle$
& $0.415^{+0.056}_{-0.055}$
& $+0.282^{+0.067}_{-0.051}$
& $10.0^{+7.0}_{-3.8}$\\

$\Bsone(5750)\to\Bs\gamma$
& $f_{\gamma,s}^{\perp}$
& $0.416^{+0.048}_{-0.054}$
& $+0.395^{+0.068}_{-0.044}$
& $19.3^{+7.9}_{-5.2}$\\

$\Bsone(5830)\to\Bs\gamma$
& $\chi_s\langle\bar s s\rangle$
& $0.246^{+0.059}_{-0.075}$
& $-0.0364^{+0.0418}_{-0.0634}$
& $0.375^{+1.74}_{-0.341}$\\

$\Bsone(5830)\to\Bs\gamma$
& $f_{\gamma,s}^{\perp}$
& $0.247^{+0.063}_{-0.075}$
& $+0.0105^{+0.0498}_{-0.0898}$
& $0.504^{+1.173}_{-0.447}$\\
\end{tabular}
\end{ruledtabular}
\end{table}
\end{widetext}

The \(\Bsone(5750)\) entry is a lower bottom-strange diagnostic configuration
rather than an established experimental resonance, and should therefore be
read as a prediction for the corresponding mixed-current state if it is
realized in nature.

The comparison of the two photon normalizations shows that the leading-twist
photon input is one of the dominant sources of uncertainty.  In the
conventional susceptibility-normalized scenario, the result depends on the
product $\chi_s\langle\bar s s\rangle$, which carries the uncertainty of both
the magnetic susceptibility and the strange-quark condensate.  The
lattice-normalized input $f_{\gamma,s}^{\perp}$ replaces this product by a
directly normalized leading-twist photon matrix element and gives the preferred
modern normalization used in the phenomenological discussion below.

The suppression pattern can be seen directly at the amplitude level.  With the
$f_{\gamma,s}^{\perp}$ input and after applying the mixing coefficients, the
two basis-current contributions to the normalized
$\Dsone(2460)\to\Ds\gamma$ amplitude are $-0.0715$ and $-0.205$; they add
constructively.  For $\Dsone(2536)\to\Ds\gamma$, the corresponding
contributions are $-0.0917$ and $+0.132$, so the physical amplitude is a small
difference.  This destructive interference explains the strong suppression of
$\Dsone(2536)\to\Ds\gamma$ and provides the LCSR realization of the
cancellation mechanism discussed in Refs.~\cite{Bondar:2023mua,Bondar:2025smw}.
A similar cancellation-sensitive pattern appears for the established
$\Bsone(5830)$ state.

\begin{figure}[t]
\centering
% \includegraphics[width=0.47\textwidth]{figures/mc_charm.pdf}
% \includegraphics[width=0.47\textwidth]{figures/mc_bottom.pdf}
\caption{Monte Carlo width distributions.  Left: charm-strange channels with
the $f_{\gamma,s}^{\perp}$ photon input.  Right: bottom-strange
physical-current normalization.}
\label{fig:mc}
\end{figure}

The distributions in Fig.~\ref{fig:mc} illustrate why percentile intervals are
used.  The $\Dsone(2460)$ distribution is relatively compact, whereas the
$\Dsone(2536)$ and $\Bsone(5830)$ distributions are distorted by cancellations
near the amplitude level.  In such cases a symmetric Gaussian uncertainty would
obscure the physical origin of the error.

Using the measured total width of the $\Dsone(2536)$, our preferred result
corresponds to a branching fraction of order
\begin{equation}
  {\cal B}[\Dsone(2536)\to\Ds\gamma]
  \sim 5\times10^{-4}.
  \label{eq:br-ds12536}
\end{equation}
This is consistent with the absence of a clear experimental signal for this
mode and provides a natural explanation for why it is difficult to isolate.
For $\Dsone(2460)$, the measured radiative branching fraction together with our
preferred partial width is compatible with a very narrow total-width scale.

The results also clarify why the predictions in the literature span a wide
range.  In cancellation-sensitive channels, small changes in the relative size
or sign of the two physical-current components can lead to large changes in
the partial width.  Future measurements of
$\Dsone(2536)\to\Ds\gamma$ and $\Bsone(5830)\to\Bs\gamma$, or improved upper
limits on these modes, would therefore provide direct tests of the spin-mixing
structure of strange heavy axial-vector mesons.


\section{Conclusion}
\label{sec:concl}

Radiative transitions of positive-parity strange heavy mesons provide a
sensitive test of the spin structure of the physical axial-vector states.  In
this work we studied the decays
$\Dsone(2460)\to\Ds\gamma$, $\Dsone(2536)\to\Ds\gamma$,
$\Bsone(5750)\to\Bs\gamma$, and $\Bsone(5830)\to\Bs\gamma$
within light-cone QCD sum rules.  The calculation was performed for two basis
axial currents, and the resulting amplitudes were combined into the physical
$1^+$ states using the same mixing convention as in the hadronic residues.  We
kept the strange-quark-mass terms and compared the conventional
susceptibility normalization of the photon distribution amplitude with the
lattice-normalized leading-twist photon input.

With the lattice-normalized photon input, our preferred charm-strange results
are
\begin{align}
  \Gamma[\Dsone(2460)\to\Ds\gamma]
  &=
  32.7^{+6.7}_{-6.2}~\keV,
  \\
  \Gamma[\Dsone(2536)\to\Ds\gamma]
  &=
  0.425^{+1.352}_{-0.390}~\keV .
\end{align}
The two widths are therefore strongly hierarchical.  The
$\Dsone(2460)\to\Ds\gamma$ width is at the tens-of-keV level, whereas
$\Dsone(2536)\to\Ds\gamma$ is suppressed to the sub-keV level.  This
suppression is traced to destructive interference between the two
physical-current components of the $\Dsone(2536)$ amplitude.  In contrast, the
corresponding components in the $\Dsone(2460)$ amplitude add constructively.
Thus the light-cone sum-rule calculation gives an independent realization of
the cancellation mechanism discussed in Refs.~\cite{Bondar:2023mua,
Bondar:2025smw}.

The result is consistent with the present experimental situation.  The
radiative decay $\Dsone(2460)\to\Ds\gamma$ has been observed, whereas no clear
signal exists for $\Dsone(2536)\to\Ds\gamma$.  Using the measured total width
of the $\Dsone(2536)$, our preferred partial width corresponds to a branching
fraction of order $5\times10^{-4}$, providing a natural explanation for why
this channel is difficult to isolate experimentally.

In the bottom-strange sector, the preferred prediction for the established
$\Bsone(5830)$ state is
\begin{equation}
  \Gamma[\Bsone(5830)\to\Bs\gamma]
  =
  0.504^{+1.173}_{-0.447}~\keV .
\end{equation}
For the lower \(\Bsone(5750)\) diagnostic configuration we find a larger
radiative width.  Since direct radiative information on bottom-strange axial
mesons is not yet available, these results provide benchmarks for future
measurements.

The comparison with earlier theoretical work shows that these radiative
transitions remain phenomenologically important.  Existing predictions span a
wide range, especially for cancellation-sensitive channels.  Our analysis
indicates that the hierarchy between the two axial-vector states is controlled
not only by phase space, but also by the relative sign and size of the two
mixed-current contributions.  Future measurements of
$\Dsone(2536)\to\Ds\gamma$ and $\Bsone(5830)\to\Bs\gamma$, or improved upper
limits on these modes, would therefore provide direct tests of the spin-mixing
structure of strange heavy axial mesons and of the radiative-transition
dynamics described here.


\appendix

\section{Photon distribution amplitudes and auxiliary functions}
\label{app:sumrule-expressions}

For the nonperturbative photon matrix elements we use the photon distribution
amplitudes (DAs) in the notation of Rohrwild~\cite{Rohrwild:2007yt}, which is
based on the BBK classification~\cite{Ball:2002ps}.  We collect here the
functions and auxiliary combinations entering the sum rules of
Sec.~\ref{sec:framework}.  The variable $u$ denotes the quark momentum
fraction, $\bar u=1-u$, and
$\alpha_q+\alpha_{\bar q}+\alpha_g=1$ for the three-particle DAs.

The leading-twist two-particle DA is taken in the asymptotic form
\begin{equation}
  \phi_\gamma(u)=6u\bar u .
  \label{eq:app-phi}
\end{equation}
The twist-3 vector DA is
\begin{align}
  \psi^{(v)}(u)
  &=
  -20u\bar u(2u-1)
  +\frac{15}{16}\left(\omega_\gamma^A-3\omega_\gamma^V\right)
  u\bar u(2u-1)\left[7(2u-1)^2-3\right].
  \label{eq:app-psiv}
\end{align}
For the twist-4 two-particle DAs we use
\begin{align}
  {\cal A}(u)
  &=
  40u\bar u\left(3\kappa-\kappa^+ +1\right)
  +8\left(\zeta_2^+-3\zeta_2\right)
  \Big[
  u\bar u(2+13u\bar u)
  \notag\\
  &\hspace{2.2cm}
  +u^3(10-15u+6u^2)\ln u
  +\bar u^3(10-15\bar u+6\bar u^2)\ln\bar u
  \Big],
  \label{eq:app-A}\\
  {\cal B}(u)
  &=
  40(1+3\kappa^+)
  \int_0^u d\alpha\,(u-\alpha)
  \left[-\frac12+\frac32(2\alpha-1)^2\right].
  \label{eq:app-B}
\end{align}

The twist-4 three-particle DAs used in the numerical analysis are
\begin{align}
  {\cal S}(\alpha_i)
  &=
  30\alpha_g^2
  \Big[
  (\kappa+\kappa^+)(1-\alpha_g)
  +(\zeta_1+\zeta_1^+)(1-\alpha_g)(1-2\alpha_g)
  \notag\\
  &\hspace{2.2cm}
  +\zeta_2\{3(\alpha_{\bar q}-\alpha_q)^2-\alpha_g(1-\alpha_g)\}
  \Big],
  \label{eq:app-S}\\
  \widetilde{\cal S}(\alpha_i)
  &=
  -30\alpha_g^2
  \Big[
  (\kappa-\kappa^+)(1-\alpha_g)
  +(\zeta_1-\zeta_1^+)(1-\alpha_g)(1-2\alpha_g)
  \notag\\
  &\hspace{2.2cm}
  +\zeta_2\{3(\alpha_{\bar q}-\alpha_q)^2-\alpha_g(1-\alpha_g)\}
  \Big],
  \label{eq:app-Stilde}\\
  {\cal T}_1(\alpha_i)
  &=
  -120(3\zeta_2+\zeta_2^+)
  (\alpha_q-\alpha_{\bar q})\alpha_{\bar q}\alpha_q\alpha_g,
  \label{eq:app-T1}\\
  {\cal T}_2(\alpha_i)
  &=
  30\alpha_g^2(\alpha_q-\alpha_{\bar q})
  \Big[
  \kappa-\kappa^+
  +(\zeta_1-\zeta_1^+)(1-2\alpha_g)
  +\zeta_2(3-4\alpha_g)
  \Big],
  \label{eq:app-T2}\\
  {\cal T}_3(\alpha_i)
  &=
  -120(3\zeta_2-\zeta_2^+)
  (\alpha_q-\alpha_{\bar q})\alpha_{\bar q}\alpha_q\alpha_g,
  \label{eq:app-T3}\\
  {\cal T}_4(\alpha_i)
  &=
  30\alpha_g^2(\alpha_q-\alpha_{\bar q})
  \Big[
  \kappa+\kappa^+
  +(\zeta_1+\zeta_1^+)(1-2\alpha_g)
  +\zeta_2(3-4\alpha_g)
  \Big].
  \label{eq:app-T4}
\end{align}

The two-particle functions entering the Borel-transformed amplitudes are used
through
\begin{align}
  \Psi^{(v)}(u_0)
  &=
  \int_0^{u_0}du\,\psi^{(v)}(u),
  \label{eq:app-Psiv-int}\\
  {\cal B}(u_0)
  &=
  -8\int_0^{u_0}du\,(u_0-u)h_\gamma(u),
  \label{eq:app-B-dictionary}
\end{align}
where Eq.~\eqref{eq:app-B-dictionary} is the notation dictionary connecting
the Rohrwild/BBK convention to the form used in the single-current benchmark
of Ref.~\cite{Colangelo:2005hv}.

For the three-particle contribution we define
\begin{align}
  {\cal I}_{F}(u_0)
  &=
  \int_0^{1-u_0}dv
  \int_0^{u_0/(1-v)}d\alpha_g\,
  F\!\left(
  u_0-(1-v)\alpha_g,\,
  1-u_0-v\alpha_g,\,
  \alpha_g,\,
  v
  \right)
  \notag\\
  &\quad
  +\int_{1-u_0}^{1}dv
  \int_0^{(1-u_0)/v}d\alpha_g\,
  F\!\left(
  u_0-(1-v)\alpha_g,\,
  1-u_0-v\alpha_g,\,
  \alpha_g,\,
  v
  \right),
  \label{eq:app-three-particle-integral}
\end{align}
with
\begin{align}
  F
  &=
  {\cal S}
  +\widetilde{\cal S}
  -{\cal T}_1
  -{\cal T}_2
  +{\cal T}_3
  +{\cal T}_4
  +2v\left(-{\cal S}-{\cal T}_3+{\cal T}_2\right).
  \label{eq:app-F-combination}
\end{align}
All three-particle DAs in Eq.~\eqref{eq:app-F-combination} are evaluated at
\((\alpha_q,\alpha_{\bar q},\alpha_g)\), with the arguments specified in
Eq.~\eqref{eq:app-three-particle-integral}.  In the numerical analysis we use
the symmetric point \(u_0=1/2\).

\begin{table}[t]
\caption{Photon-DA parameters used at $\mu=1~\GeV$ in Rohrwild notation.}
\label{tab:app-photon-params}
\begin{ruledtabular}
\begin{tabular}{lll}
Parameter & Value & Comment\\
\hline
$\chi_s$ & $3.15\pm0.30~\GeV^{-2}$ & legacy susceptibility input\\
$f_{3\gamma}$ & $-(4.0\pm2.0)\times10^{-3}~\GeV^2$ & twist-3 coupling\\
$\omega_\gamma^A$ & $-2.1\pm1.0$ & twist-3 shape\\
$\omega_\gamma^V$ & $3.8\pm1.8$ & twist-3 shape\\
$\kappa$ & $0.15$ & twist-4 shape\\
$\kappa^+$ & $-0.05$ & twist-4 shape\\
$\zeta_1$ & $0.40$ & twist-4 shape\\
$\zeta_1^+$ & $0$ & twist-4 shape\\
$\zeta_2$ & $0.30$ & twist-4 shape\\
$\zeta_2^+$ & $0$ & twist-4 shape\\
\end{tabular}
\end{ruledtabular}
\end{table}

In the preferred modern scenario we replace the legacy normalization
$\chi_s\langle\bar s s\rangle$ of the leading-twist matrix element by the
lattice-normalized input $f_{\gamma,s}^{\perp}$.  The DA shapes in this
appendix are otherwise kept fixed in the same Rohrwild/BBK convention.



\section{Basis-current invariant amplitudes}
\label{app:basis-amplitudes}

We write the Borel-transformed invariant amplitudes in the form
\begin{equation}
  \widehat\Pi_K(M^2,s_0)
  =
  \widehat\Pi_K^{\rm pert}
  +\widehat\Pi_K^{(2p)}
  +\widehat\Pi_K^{(3p)},
  \qquad K=A,B ,
  \label{eq:app-pi-decomp}
\end{equation}
where the three terms denote the perturbative photon emission, two-particle
photon-DA contribution, and three-particle photon-DA contribution,
respectively.  At the symmetric Borel point used in the numerical analysis,
\[
  M_1^2=M_2^2=2M^2,\qquad u_0=\frac12 ,
\]
the perturbative parts are written as
\begin{equation}
  \widehat\Pi_K^{\rm pert}
  =
  \int_{(m_Q+m_s)^2}^{s_0}ds\,e^{-s/M^2}\rho_K(s),
  \qquad K=A,B .
  \label{eq:app-pert}
\end{equation}





\begin{acknowledgments}
To be added.
\end{acknowledgments}

\bibliographystyle{apsrev4-2}
\bibliography{../../../shared/bibliography/all}

\end{document}
